\documentclass[11pt,a4paper]{article}
\usepackage[margin=24mm,headheight=14pt]{geometry}
\usepackage{amsmath,amssymb,amsthm,mathtools,booktabs}
\usepackage{lmodern,microtype,fancyhdr,fontspec,fancyvrb}
\usepackage[hidelinks]{hyperref}
\newfontfamily\leanappendixfont{DejaVuSansMono.ttf}
\newtheorem{theorem}{Theorem}
\newtheorem{lemma}[theorem]{Lemma}
\newtheorem{proposition}[theorem]{Proposition}
\newcommand{\code}[1]{\texttt{#1}}
\pagestyle{fancy}\fancyhf{}
\fancyhead[L]{\small Consecutive integers with equal Collatz height}
\fancyhead[R]{\small Lucas Valbuena}\fancyfoot[C]{\thepage}
\renewcommand{\headrulewidth}{0.3pt}
\setlength{\parindent}{1.1em}\setlength{\parskip}{3pt}
\setlength{\abovedisplayskip}{7pt plus 2pt minus 2pt}
\setlength{\belowdisplayskip}{7pt plus 2pt minus 2pt}
\hypersetup{pdftitle={Arbitrarily long runs of consecutive integers with equal Collatz height},pdfauthor={Lucas Valbuena},pdfsubject={Finite-pattern synchronization and equal finite first hitting times; Lean-checked proofs}}
\begin{document}\thispagestyle{plain}
\begin{center}
{\LARGE\bfseries Arbitrarily long runs of consecutive integers\\[3pt]
with equal Collatz height\par}
\vspace{8pt}{\large Lucas Valbuena\par}
\vspace{4pt}{\normalsize 28 September 2026}
\end{center}
\begin{abstract}
For every prescribed length and lower bound, we construct a larger
interval of consecutive positive integers that first reach one after
the same number of shortcut Collatz steps, with equal odd-step counts.
More generally, every finite pattern of nonnegative offsets has
arbitrarily large translates with these properties. The proof first
coalesces the pattern on a dyadic progression, then specializes its
common endpoint to a power of two. Equal first hitting times also
follow for the unshortened map. The principal statements and their
supporting lemmas are formalized in Lean~4. These results do not
establish convergence from an arbitrary prescribed starting integer.
\end{abstract}

\section{Heights and the main result}
Let $\mathbb N$ include zero. For positive integers, put
\[
T(x)=\begin{cases}x/2,&x\text{ even},\\(3x+1)/2,&x\text{ odd}.
\end{cases}
\]
We use \emph{height} for the first hitting time
$h(x)=\min\{k\in\mathbb N:T^k(x)=1\}$, taking $h(x)=\infty$
when that set is empty. In particular, $h(1)=0$.
Write $o_k(x)$ for the number of odd inputs among
$x,T(x),\ldots,T^{k-1}(x)$, with $o_0(x)=0$.
This is the height convention used by Garner~\cite{Garner1985},
as distinguished from other conventions in
LaDue~\cite[p.~156]{LaDue2018}.

Garner's 1985 run-length conjecture~\cite{Garner1985}, as recorded
by Lagarias~\cite[entry~72, p.~29]{Lagarias2011}, asserts that runs
of consecutive integers with equal height can be arbitrarily long.
Theorem~\ref{thm:equal-heights} proves this assertion with equal odd
counts as well, and with the runs occurring above every prescribed
lower bound. A separate conjecture of Garner concerned a proposed
characterization of equal-height consecutive pairs; the counterexamples
of Elia and Tucker~\cite{EliaTucker2015} concern that characterization,
not the run-length assertion. Related studies include Gao's computations
of long runs~\cite{Gao1993} and LaDue's results on
clusters~\cite{LaDue2018}. We give a self-contained construction and
its Lean formalization, without claiming historical priority.


\begin{theorem}\label{thm:equal-heights}
For every integer $\ell\ge1$ and every $B\in\mathbb N$, there are
an integer $x>\max\{1,B\}$ and $H,R\in\mathbb N$ such that
\[
h(x+i)=H,\qquad o_H(x+i)=R\qquad(0\le i<\ell).
\]
Thus arbitrarily long runs of consecutive integers have the same
finite Collatz height, and each prescribed run length occurs above
every lower bound.
\end{theorem}

For example, $28,29,30$ all reach $20$ after seven steps and first
reach one after thirteen steps, with five odd inputs each.
We prove a stronger finite-pattern statement: any finite set
$F\subseteq\mathbb N$ may replace $\{0,\ldots,\ell-1\}$ in the theorem.
The translating integer $x$ is chosen by the construction; no
convergence assertion for an arbitrary fixed $x$ is assumed or deduced.
The power-of-two endpoint argument uses the fact that two generates
the units modulo powers of three, also used by
Monks~\cite[Lemma~3.4]{Monks2006} in back tracing.

\section{Coalescence on a dyadic progression}

We first record the dependence of a finite trajectory on its starting
residue class.

\begin{lemma}\label{lem:finite-dyadic-prefix}
For $c\in\mathbb Z_{>0}$, $k,t\in\mathbb N$, and $0\le j\le k$,
\begin{align}
T^j(c+2^kt)
 &=T^j(c)+2^{k-j}3^{o_j(c)}t,
 \label{eq:finite-dyadic-prefix}\\
o_j(c+2^kt)&=o_j(c).
 \label{eq:finite-dyadic-count}
\end{align}
\end{lemma}
\begin{proof}
The assertions hold at $j=0$. If $j<k$, the difference in
\eqref{eq:finite-dyadic-prefix} is even, so the two current integers
have the same parity. An even step divides the difference by two;
an odd step multiplies it by $3/2$ and adds one to both odd counts.
This proves the induction.
\end{proof}

In particular, the first $k$ parities are constant on each positive
residue class modulo $2^k$. If a current value is $A+3^rt$, either
parity can be prescribed by restricting $t$ to $\varepsilon+2u$,
where $\varepsilon\in\{0,1\}$. After the prescribed actual step,
the coefficient of $u$ is $3^r$ or $3^{r+1}$ according as that
step is even or odd. Thus this prescription can be repeated, and
each additional step doubles the parameter modulus.

\begin{lemma}\label{lem:finite-pair}
For every $n\in\mathbb N$, there are $k\in\mathbb N$ and
$a\in\mathbb Z_{>0}$
such that, for every $t\in\mathbb N$,
\begin{align*}
T^k(a+2^kt)&=T^k(a+2^kt+n),\\
o_k(a+2^kt)&=o_k(a+2^kt+n).
\end{align*}
\end{lemma}
\begin{proof}
We use strong induction on $n$. For $n=0$, take $k=0$ and $a=1$.
For $n=1$, the two trajectories
\[
\begin{aligned}
4+8t&\longmapsto2+4t\longmapsto1+2t\longmapsto2+3t,\\
5+8t&\longmapsto8+12t\longmapsto4+6t\longmapsto2+3t
\end{aligned}
\]
each have one odd step. This gives $k=3$ and $a=4$.

If $n=2m>0$, use the progression $a+2^kt$ for $m$ supplied
by induction. The starts $2a+2^{k+1}t$ and
$2a+2^{k+1}t+2m$ take one even step to that pair. They therefore
coalesce after $k+1$ steps, with equal odd counts.

It remains to treat odd $n>1$. Choose $m,c,d$ by
\[
\begin{array}{c|ccc}
n\bmod4&m&c&d\\ \hline
1&(3n+1)/4&1&2\\
3&(3n-1)/4&2&1
\end{array}
\]
so that $0<m<n$ and $3n+d=4m+c$.
Let $a+2^kt$ be the progression for $m$.
Because $2^k$ is invertible modulo three, some
$q\in\{0,1,2\}$ makes $4(a+2^kq)-c$ divisible by three.
The integer
\[
a'=\frac{4(a+2^kq)-c}{3}
\]
is positive. For $X=a'+2^{k+2}t$ and
$Y=a+2^k(q+3t)$, we have
\[
3X+c=4Y,\qquad 3(X+n)+d=4(Y+m).
\]
If $3z+1=4w$, then $z\equiv1\pmod4$ and its next two
steps are odd, even, with endpoint $w$.
If $3z+2=4w$, then $z\equiv2\pmod4$ and its next two
steps are even, odd, again with endpoint $w$.
Hence $X$ and $X+n$ reach $Y$ and $Y+m$ after two steps,
each having used one odd step. The induction hypothesis at parameter
$q+3t$ finishes the construction.
\end{proof}

\subsection{Synchronizing affine families}

To add an offset to a collection that already coalesces, we must
synchronize its current trajectory with the common trajectory of the
collection. Their coefficients in the remaining parameter can be
different powers of three.

\begin{lemma}\label{lem:finite-affine-synchronization}
Let $A,B\in\mathbb Z_{>0}$ and $r,s\in\mathbb N$.
There exist $\ell,t_0,Q\in\mathbb N$ and
$C\in\mathbb Z_{>0}$ such that,
for every $v\in\mathbb N$,
\begin{equation}\label{eq:finite-affine-synchronization}
T^\ell\bigl(A+3^r(t_0+2^\ell v)\bigr)
=T^\ell\bigl(B+3^s(t_0+2^\ell v)\bigr)
=C+3^Qv.
\end{equation}
The two length-$\ell$ trajectories have fixed odd counts $u,w$
satisfying $r+u=s+w=Q$.
\end{lemma}
\begin{proof}
Exchange the families if necessary so that $r\le s$.
Initially their current values satisfy
\[
Y=3^dX+D,\qquad d=s-r,\qquad D=B-3^dA\in\mathbb Z.
\]
At every subsequent stage the coefficient of the remaining parameter
in $X$ is a power of three and is therefore odd. We may prescribe its next
parity, as above. The displayed relation determines the parity of $Y$.
When $d>0$, the following steps reduce the relation to one with a
smaller exponent.

If $D\ne0$ is even, prescribe $X$ even. Both values take an even
step, and
\[
Y'=3^dX'+D/2.
\]
Repeating this finitely many times makes $D$ odd. This also applies
to negative $D$, by dividing out the power of two in $|D|$.
If $D=0$, prescribe $X$ odd. Both values take an odd step, and
\[
Y'=3^dX'+\frac{1-3^d}{2}.
\]
The new constant is nonzero because $d>0$, so the preceding
halving procedure makes it odd.
Once $D$ is odd, prescribe $X$ odd. Then $Y$ is even, and
\begin{equation}\label{eq:finite-ratio-reduction}
X'=\frac{3X+1}{2},\qquad Y'=\frac Y2,
\qquad
Y'=3^{d-1}X'+\frac{D-3^{d-1}}2.
\end{equation}
The last constant is an integer and the exponent has decreased.
There are finitely many steps at each fixed positive $d$ and finitely
many decreases of $d$, so eventually $d=0$.

The current values now have the same coefficient in the remaining
parameter and differ by a constant integer $D$. If $D=0$, they
are equal already. Otherwise, apply Lemma~\ref{lem:finite-pair}
with gap $|D|$ to the smaller current value. Suppose its required
progression is $a+2^kz$. The smaller value has the form
$E+3^qu$, so the congruence
\[
E+3^qu_0\equiv a\pmod{2^k}
\]
has a solution. Choose a nonnegative representative $u_0$ large
enough that $E+3^qu_0\ge a$. On the parameter class
$u=u_0+2^kv$, the smaller value belongs to the positive tail of
the required progression for every $v\ge0$.
The next $k$ steps therefore make the two values equal.

Each parameter restriction has modulus two for each added actual
step. Their composition has the form $t=t_0+2^\ell v$, where
$\ell$ is the common total number of steps. All bases and current
values remain positive. By Lemma~\ref{lem:finite-dyadic-prefix},
the resulting odd counts $u,w$ are independent of $v$, and the
two final affine coefficients are $3^{r+u}$ and $3^{s+w}$.
Equality at $v=0$ and $v=1$ forces these coefficients to agree.
This proves \eqref{eq:finite-affine-synchronization} and
$r+u=s+w=Q$.
\end{proof}

\begin{proposition}\label{prop:finite-pattern-coalescence}
For every finite $F\subset\mathbb N$, there are $K,R\in\mathbb N$
and $a,b\in\mathbb Z_{>0}$ such that, for every $n\in F$ and $t\in\mathbb N$,
\begin{equation}\label{eq:finite-pattern-coalescence}
T^K(a+2^Kt+n)=b+3^Rt,\qquad
o_K(a+2^Kt+n)=R.
\end{equation}
\end{proposition}
\begin{proof}
Enlarge $F$ to include zero, and enumerate its offsets starting with
zero. For the initial singleton take $K=R=0$ and $a=b=1$.
Suppose the assertion holds for the offsets already included.
Their common current value on $a+2^Kt$ is $U(t)=b+3^Rt$.
For the next offset $n$, Lemma~\ref{lem:finite-dyadic-prefix}
gives
\[
T^K(a+2^Kt+n)=d+3^St,
\qquad d=T^K(a+n)>0,\quad S=o_K(a+n).
\]
Apply Lemma~\ref{lem:finite-affine-synchronization} to these two
affine families. On a class $t=t_0+2^\ell v$, they have the
same image after $\ell$ additional steps. Every previously included
trajectory starts those steps at exactly $U(t)$, so all previous
equalities are preserved.

The new starting progression is
\[
(a+2^Kt_0)+2^{K+\ell}v.
\]
Its modulus agrees with the new common time $K+\ell$.
If the two added trajectories use $u,w$ odd steps, the last
assertion of Lemma~\ref{lem:finite-affine-synchronization} gives
$R+u=S+w$. Thus all total odd counts agree as well.
The new endpoint has the required form $b'+3^{R+u}v$.
Finite induction proves the proposition.
\end{proof}

\section{Equal heights by endpoint specialization}

\begin{lemma}\label{lem:finite-power-two-residue}
Let $b\in\mathbb Z_{>0}$ with $3\nmid b$, and let $R\in\mathbb N$.
There are arbitrarily large $t\in\mathbb N$ for which
$b+3^Rt$ is a power of two.
\end{lemma}
\begin{proof}
We give the residue argument explicitly. For $m\ge1$,
\begin{equation}\label{eq:finite-ternary-lift}
4^{3^{m-1}}\equiv1+3^m\pmod{3^{m+1}}.
\end{equation}
The case $m=1$ is immediate. Cubing a number of the form
$1+3^m(1+3c)$ proves the next case, since the terms beyond the
linear term are divisible by $3^{m+2}$.

Powers of two represent both units modulo three. Suppose
$2^s\equiv b\pmod{3^m}$. Multiplying $2^s$ by
$(4^{3^{m-1}})^j$, for $j=0,1,2$, changes its residue modulo
$3^{m+1}$ by $j2^s3^m$, by \eqref{eq:finite-ternary-lift}.
Since $3\nmid2^s$, these are the three possible lifts of its
residue modulo $3^m$. One therefore represents $b$ modulo
$3^{m+1}$. Induction gives $2^s\equiv b\pmod{3^R}$
for every $R\ge1$; the case $R=0$ is immediate.

The same identity gives
$2^{2\cdot3^R}\equiv1\pmod{3^R}$.
We may consequently increase $s$ by arbitrary multiples of
$2\cdot3^R$ while preserving the residue. For all sufficiently
large such exponents,
\[
t=\frac{2^s-b}{3^R}
\]
is a nonnegative integer, and these values of $t$ are unbounded.
\end{proof}

\begin{proof}[Proof of Theorem~\ref{thm:equal-heights}]
We prove the stronger assertion for any finite set $F\subseteq\mathbb N$
of offsets, which includes $F=\{0,\ldots,\ell-1\}$.
Apply Proposition~\ref{prop:finite-pattern-coalescence} to
$F\cup\{0,1\}$, obtaining $K,R,a,b$.
We first show $3\nmid b$. An odd step has output congruent
to two modulo three. Thus a current value divisible by three has
only an even predecessor, which is also divisible by three.
If $3\mid b$, every step in any length-$K$ trajectory ending
at $b$ would be even, and its initial value would equal $2^Kb$.
The distinct starts $a$ and $a+1$ both end at $b$, a
contradiction.

By Lemma~\ref{lem:finite-power-two-residue}, choose $t\ge1$
as large as necessary, with $b+3^Rt=2^s$, and put
$x=a+2^Kt$. The common endpoint after $K$ steps is $2^s$.
After $s$ further even steps, all trajectories reach one, with
total odd count $R$. Taking $t$ large ensures
$x>\max\{1,B\}$.

It remains to check that this is the first visit to one.
For $j<K$, Lemma~\ref{lem:finite-dyadic-prefix} gives
\[
T^j(x+n)=T^j(a+n)+2^{K-j}3^{o_j(a+n)}t>1.
\]
For $j=K+q<K+s$, the value is $2^{s-q}>1$.
Thus $h(x+n)=H=K+s$ and $o_H(x+n)=R$ for every $n\in F$.
Taking $F=\{0,\ldots,\ell-1\}$ gives the theorem.
\end{proof}

For the unshortened map that sends odd $x$ to $3x+1$, each
of the $R$ odd steps above expands to an odd step followed by an
even step. The corresponding first hitting time is therefore
$H+R$, also common to the whole pattern. The inserted intermediate
values are greater than one, so first arrival is preserved.


\section{Formal verification}
The independent artifact contains nine Lean modules, with 66 public
theorems and lemmas and eight definitions. A fresh build and axiom
audit checks all 74 declarations using Lean~4.27.0 and mathlib commit
\begin{center}\scriptsize
\code{a3a10db0e9d66acbebf76c5e6a135066525ac900}.
\end{center}
Only \code{propext}, \code{Classical.choice}, and \code{Quot.sound}
are permitted. No proof placeholder, added mathematical axiom, or
numerical search result is used as proof. The entry point is
\path{formal/FinitePatternConvergence.lean}; its two theorems state
the arbitrary finite-pattern result and the consecutive-run corollary.
Both explicitly require every earlier iterate to exceed one.
Appendix~\ref{app:lean} reproduces the definitions and theorem types.
The elementary conversion to the unshortened map is a prose consequence;
the Lean endpoints use $T$ as defined above.

The \path{equal-collatz-heights/} directory contains the source, PDF,
claim map, pinned verifier, and retained build report. The repository is
\begin{center}\scriptsize
\url{https://github.com/x1xhlol/collatz-two-coordinate-obstruction}.
\end{center}
From that directory, run \code{python3 verify.py} with
\code{--mathlib-root} pointing to the pinned checkout and
\code{--report} to an unused output path. The accompanying README gives
the complete setup. Kernel acceptance verifies the formal proposition;
matching that proposition to the mathematical statement is a separate
reading task addressed by the appendix and claim map.

\paragraph{Acknowledgments and use of AI.}
The formalization uses Lean and mathlib. OpenAI Codex was used for
mathematical exploration, proof development, verification scripts,
and manuscript preparation. The claims are supported by the
accompanying proof artifacts; this disclosure implies no endorsement
by authors of the cited work.

\begin{thebibliography}{9}\small\raggedright
\bibitem{Garner1985} L. E. Garner.
\emph{On heights in the Collatz $3n+1$ problem}.
Discrete Mathematics \textbf{55}(1), 57--64, 1985.
\url{https://doi.org/10.1016/S0012-365X(85)80020-0}.

\bibitem{Gao1993} G.-G. Gao.
\emph{On consecutive numbers of the same height in the Collatz problem}.
Discrete Mathematics \textbf{112}(1--3), 261--267, 1993.
\url{https://doi.org/10.1016/0012-365X(93)90240-T}.

\bibitem{Monks2006} K. M. Monks.
\emph{The sufficiency of arithmetic progressions for the $3x+1$ Conjecture}.
Proceedings of the American Mathematical Society \textbf{134}(10),
2861--2872, 2006.
\url{https://doi.org/10.1090/S0002-9939-06-08567-4}.
Author version:
\url{https://monks.scranton.edu/files/pubs/SufficiencyRev4.pdf}.

\bibitem{LaDue2018} M. D. LaDue.
\emph{Clusters of integers with equal total stopping times in the $3x+1$ problem}.
The Fibonacci Quarterly \textbf{56}(2), 156--162, 2018.
\url{https://www.fq.math.ca/Papers1/56-2/ladue121517.pdf}.


\bibitem{Lagarias2011} J. C. Lagarias.
\emph{The $3x+1$ Problem: An Annotated Bibliography (1963--1999)}.
arXiv:math/0309224v13, 2011, entry~72, p.~29.
\url{https://arxiv.org/pdf/math/0309224v13}.

\bibitem{EliaTucker2015} M. Elia and A. Tucker.
\emph{Consecutive Integers and the Collatz Conjecture}.
arXiv:1511.09141, 2015.
\url{https://arxiv.org/abs/1511.09141}.

\end{thebibliography}
\clearpage\appendix
\section{Exact Lean declarations}\label{app:lean}
The following declarations use namespace
\path{CollatzPositiveProgression}.
Natural numbers include zero. The constructed starts are positive.
Here \code{oddCount k x} counts odd inputs in the first $k$ steps.
The first-hitting-time theorem requires every earlier value to
exceed one.

% Source SHA256 PositiveProgressionCoalescence.lean: 9f503a2ea7ca75c8dcdfe8b7de95618eb526b5ee5bea8c5bf69bb8c35396d58f
\noindent\begin{minipage}{\linewidth}
\paragraph{The map, its iterates, and its odd-step count.}\mbox{}\par
\noindent{\small\path{formal/PositiveProgressionCoalescence.lean}, lines 9--18.}\par
\begin{Verbatim}[formatcom=\leanappendixfont,fontsize=\scriptsize,numbers=left,firstnumber=9,numbersep=5pt,xleftmargin=18pt]
def step (x : ℕ) : ℕ :=
  if x % 2 = 0 then x / 2 else (3 * x + 1) / 2

def iterate : ℕ → ℕ → ℕ
  | 0, x => x
  | k + 1, x => iterate k (step x)

def oddCount : ℕ → ℕ → ℕ
  | 0, _ => 0
  | k + 1, x => x % 2 + oddCount k (step x)
\end{Verbatim}
\end{minipage}\par\medskip

% Source SHA256 FinitePatternCoalescence.lean: 86208b374e1df4c1ecbefcd33270853b3b0a3349fdeb0aad854ae22939729087
\noindent\begin{minipage}{\linewidth}
\paragraph{Coalescence on a common dyadic progression.}\mbox{}\par
\noindent{\small\path{formal/FinitePatternCoalescence.lean}, lines 50--53.}\par
\begin{Verbatim}[formatcom=\leanappendixfont,fontsize=\scriptsize,numbers=left,firstnumber=50,numbersep=5pt,xleftmargin=18pt]
theorem every_finite_pattern_coalesces (F : Finset ℕ) :
    ∃ K R a b : ℕ, 0 < a ∧ 0 < b ∧ ∀ n ∈ F, ∀ t : ℕ,
      iterate K (a + 2 ^ K * t + n) = b + 3 ^ R * t ∧
      oddCount K (a + 2 ^ K * t + n) = R
\end{Verbatim}
\end{minipage}\par\medskip

% Source SHA256 FinitePatternConvergence.lean: 74ccd96f004198fbd04c6a5283022cb78414fd0351f49bc4c0ad61fc1c30ffb5
\noindent\begin{minipage}{\linewidth}
\paragraph{Equal finite first hitting times for any finite pattern.}\mbox{}\par
\noindent{\small\path{formal/FinitePatternConvergence.lean}, lines 8--12.}\par
\begin{Verbatim}[formatcom=\leanappendixfont,fontsize=\scriptsize,numbers=left,firstnumber=8,numbersep=5pt,xleftmargin=18pt]
theorem every_finite_pattern_has_equal_first_hitting_times
    (F : Finset ℕ) (lower : ℕ) :
    ∃ x H R : ℕ, lower < x ∧ 1 < x ∧ ∀ n ∈ F,
      iterate H (x + n) = 1 ∧ oddCount H (x + n) = R ∧
      ∀ j < H, 1 < iterate j (x + n)
\end{Verbatim}
\end{minipage}\par\medskip

% Source SHA256 FinitePatternConvergence.lean: 74ccd96f004198fbd04c6a5283022cb78414fd0351f49bc4c0ad61fc1c30ffb5
\noindent\begin{minipage}{\linewidth}
\paragraph{Consecutive runs of every prescribed length.}\mbox{}\par
\noindent{\small\path{formal/FinitePatternConvergence.lean}, lines 32--36.}\par
\begin{Verbatim}[formatcom=\leanappendixfont,fontsize=\scriptsize,numbers=left,firstnumber=32,numbersep=5pt,xleftmargin=18pt]
theorem arbitrarily_long_consecutive_equal_first_hitting_times
    (length lower : ℕ) :
    ∃ x H R : ℕ, lower < x ∧ 1 < x ∧ ∀ n < length,
      iterate H (x + n) = 1 ∧ oddCount H (x + n) = R ∧
      ∀ j < H, 1 < iterate j (x + n)
\end{Verbatim}
\end{minipage}\par\medskip

\end{document}
